\documentclass[12pt,a4paper]{article}

\usepackage[T2A]{fontenc}
\usepackage[utf8]{inputenc}
\usepackage[russian]{babel}
\usepackage{amsmath,amssymb}
\usepackage{enumitem}
\usepackage{array}
\usepackage{longtable}
\usepackage{listings}
\usepackage{xcolor}
\usepackage{geometry}
\usepackage{tikz}
\usetikzlibrary{arrows.meta,positioning}

\geometry{left=2.2cm,right=2.2cm,top=2cm,bottom=2cm}

\lstset{
    basicstyle=\ttfamily\small,
    keywordstyle=\color{blue},
    commentstyle=\color{gray},
    stringstyle=\color{teal},
    showstringspaces=false,
    frame=single,
    breaklines=true
}

\title{Билет №65 \\ \small Виды процессов и управления ими в современных ОС. Представление процессов, их контексты, иерархии порождения, состояния и взаимодействие. Многозадачный (многопрограммный) режим работы. Команды управления процессами. Средства взаимодействия процессов. (ИТМО)}
\author{Сериков Владислав Александрович}
\date{16 мая 2026}

\begin{document}

\maketitle
\tableofcontents
\newpage

\section{Основные понятия}

\subsection{Программа, процесс, поток}

\textbf{Программа} --- это пассивный объект: файл или набор файлов, содержащих исполняемый код, данные, метаданные, таблицы символов, библиотеки и другую информацию.

\textbf{Процесс} --- это экземпляр выполняющейся программы вместе с выделенными ей ресурсами и текущим состоянием выполнения.

Процесс включает:
\begin{itemize}
    \item исполняемый код;
    \item виртуальное адресное пространство;
    \item открытые файлы и дескрипторы;
    \item значения регистров процессора;
    \item стек вызовов;
    \item кучу;
    \item информацию о правах доступа;
    \item сведения о родителе, потомках, группе процессов;
    \item состояние планирования.
\end{itemize}

\textbf{Поток} --- это минимальная единица планирования, выполняющаяся внутри процесса. Несколько потоков одного процесса обычно разделяют адресное пространство, открытые файлы и другие ресурсы процесса, но имеют собственные:
\begin{itemize}
    \item счетчик команд;
    \item регистровый контекст;
    \item стек;
    \item состояние планирования.
\end{itemize}

\subsection{Процесс как абстракция ОС}

Процесс является одной из центральных абстракций операционной системы. Он позволяет:
\begin{itemize}
    \item изолировать программы друг от друга;
    \item распределять процессорное время;
    \item контролировать доступ к памяти, файлам и устройствам;
    \item организовывать многозадачность;
    \item вести учет ресурсов;
    \item обеспечивать безопасность и разграничение прав.
\end{itemize}

С точки зрения пользователя процесс выглядит как выполняющаяся программа. С точки зрения ядра ОС процесс --- это запись в системных структурах данных, прежде всего в таблице процессов.

\section{Виды процессов}

\subsection{По назначению}

\begin{enumerate}
    \item \textbf{Пользовательские процессы} --- выполняют прикладные программы пользователей: браузеры, редакторы, компиляторы, серверы приложений.
    \item \textbf{Системные процессы} --- обслуживают работу ОС: демоны, службы, процессы управления устройствами, планировщики задач, журналирование.
    \item \textbf{Ядерные процессы или потоки ядра} --- выполняют задачи внутри ядра, например обработку отложенных операций, ввод-вывод, обслуживание файловых систем.
\end{enumerate}

\subsection{По режиму выполнения}

\begin{enumerate}
    \item \textbf{Интерактивные процессы} --- активно взаимодействуют с пользователем и требуют малой задержки отклика.

    Примеры:
    \begin{itemize}
        \item командная оболочка;
        \item текстовый редактор;
        \item графический интерфейс.
    \end{itemize}

    \item \textbf{Пакетные процессы} --- выполняют длительную работу без постоянного взаимодействия с пользователем.

    Примеры:
    \begin{itemize}
        \item компиляция большого проекта;
        \item резервное копирование;
        \item обработка данных.
    \end{itemize}

    \item \textbf{Процессы реального времени} --- должны завершать операции в строго заданные временные интервалы.

    Они делятся на:
    \begin{itemize}
        \item \textbf{жесткого реального времени}, где нарушение срока недопустимо;
        \item \textbf{мягкого реального времени}, где нарушение срока ухудшает качество, но не приводит к катастрофе.
    \end{itemize}
\end{enumerate}

\subsection{По характеру использования процессора и ввода-вывода}

\begin{enumerate}
    \item \textbf{CPU-bound процессы} --- большую часть времени используют процессор.

    Примеры:
    \begin{itemize}
        \item численные расчеты;
        \item сжатие данных;
        \item рендеринг;
        \item криптографические вычисления.
    \end{itemize}

    \item \textbf{I/O-bound процессы} --- часто ожидают операции ввода-вывода.

    Примеры:
    \begin{itemize}
        \item файловые серверы;
        \item базы данных;
        \item сетевые приложения;
        \item терминальные программы.
    \end{itemize}
\end{enumerate}

\subsection{По отношению к терминалу}

В UNIX-подобных ОС часто различают:

\begin{enumerate}
    \item \textbf{Процессы переднего плана} --- связаны с терминалом и получают ввод с клавиатуры.
    \item \textbf{Фоновые процессы} --- выполняются без монопольного владения терминалом.
    \item \textbf{Демоны} --- фоновые системные процессы, обычно не связанные с управляющим терминалом.
\end{enumerate}

Пример запуска фонового процесса:

\begin{lstlisting}[language=bash]
sleep 100 &
\end{lstlisting}

\section{Представление процесса в ОС}

\subsection{Блок управления процессом}

Для каждого процесса ядро хранит специальную структуру данных --- \textbf{PCB} (\textit{Process Control Block}), или блок управления процессом.

В Linux близкую роль играет структура \texttt{task\_struct}.

Типичный PCB содержит:

\begin{itemize}
    \item идентификатор процесса;
    \item идентификатор родительского процесса;
    \item состояние процесса;
    \item приоритет;
    \item счетчик команд;
    \item содержимое регистров;
    \item указатели на таблицы страниц;
    \item информацию о памяти;
    \item список открытых файлов;
    \item учетную информацию;
    \item сигнальные маски;
    \item сведения о группе процессов и сеансе;
    \item информацию для планировщика.
\end{itemize}

\subsection{Схематическое представление процесса}

\begin{center}
\begin{tikzpicture}[
    block/.style={draw, rectangle, rounded corners, minimum width=7cm, minimum height=0.8cm, align=center},
    node distance=0.25cm
]
\node[block] (pcb) {PCB: PID, PPID, состояние, приоритет, регистры};
\node[block, below=of pcb] (code) {Код программы};
\node[block, below=of code] (data) {Глобальные данные};
\node[block, below=of data] (heap) {Куча};
\node[block, below=of heap] (stack) {Стек};
\node[block, below=of stack] (files) {Открытые файлы и дескрипторы};
\end{tikzpicture}
\end{center}

\section{Контекст процесса}

\subsection{Понятие контекста}

\textbf{Контекст процесса} --- это вся информация, необходимая для приостановки процесса и последующего продолжения его выполнения с того же места.

Контекст можно разделить на несколько частей:

\begin{enumerate}
    \item \textbf{Пользовательский контекст}
    \begin{itemize}
        \item код;
        \item данные;
        \item куча;
        \item пользовательский стек;
        \item содержимое пользовательского адресного пространства.
    \end{itemize}

    \item \textbf{Регистровый контекст}
    \begin{itemize}
        \item счетчик команд;
        \item указатель стека;
        \item регистры общего назначения;
        \item регистры состояния процессора;
        \item регистры управления памятью.
    \end{itemize}

    \item \textbf{Системный контекст}
    \begin{itemize}
        \item PCB;
        \item таблицы страниц;
        \item открытые файлы;
        \item права доступа;
        \item сигналы;
        \item параметры планирования.
    \end{itemize}
\end{enumerate}

\subsection{Переключение контекста}

\textbf{Переключение контекста} --- это операция ядра, при которой процессор прекращает выполнение одного процесса и начинает или продолжает выполнение другого.

Переключение контекста происходит:
\begin{itemize}
    \item при истечении кванта времени;
    \item при блокировке процесса на вводе-выводе;
    \item при ожидании синхронизационного объекта;
    \item при появлении более приоритетного процесса;
    \item при системном вызове, который приводит к планированию;
    \item при прерывании.
\end{itemize}

\subsection{Алгоритм переключения контекста}

\textbf{Вход:} текущий процесс $P_i$, следующий процесс $P_j$.

\textbf{Выход:} процессор выполняет $P_j$.

\begin{enumerate}
    \item Возникает событие, требующее вмешательства ядра: таймер, системный вызов, прерывание, блокировка.
    \item Процессор переходит в режим ядра.
    \item Ядро сохраняет регистровый контекст $P_i$ в его PCB.
    \item Обновляется состояние $P_i$:
    \begin{itemize}
        \item \texttt{Running} $\rightarrow$ \texttt{Ready}, если процесс вытеснен;
        \item \texttt{Running} $\rightarrow$ \texttt{Blocked}, если процесс ожидает событие.
    \end{itemize}
    \item Планировщик выбирает следующий процесс $P_j$ из очереди готовых процессов.
    \item Ядро загружает адресное пространство $P_j$.
    \item Ядро восстанавливает регистры $P_j$ из его PCB.
    \item Состояние $P_j$ меняется на \texttt{Running}.
    \item Процессор возвращается в пользовательский режим и продолжает выполнение $P_j$.
\end{enumerate}

\subsection{Иллюстрация переключения контекста}

\begin{center}
\begin{tikzpicture}[
    proc/.style={draw, rectangle, rounded corners, minimum width=3.8cm, minimum height=1cm, align=center},
    arrow/.style={->, thick}
]
\node[proc] (p1) {Процесс $P_1$\\Running};
\node[proc, right=3cm of p1] (kernel) {Ядро ОС\\сохраняет/загружает контекст};
\node[proc, right=3cm of kernel] (p2) {Процесс $P_2$\\Running};

\draw[arrow] (p1) -- node[above]{прерывание} (kernel);
\draw[arrow] (kernel) -- node[above]{dispatch} (p2);
\end{tikzpicture}
\end{center}

Переключение контекста является накладной операцией: во время него пользовательская работа не выполняется.

\section{Иерархия порождения процессов}

\subsection{Родительские и дочерние процессы}

Во многих ОС процесс может создавать другие процессы.

\textbf{Родительский процесс} --- процесс, создавший новый процесс.

\textbf{Дочерний процесс} --- процесс, созданный другим процессом.

В UNIX-подобных системах эта модель особенно выражена:
\begin{itemize}
    \item каждый процесс имеет PID;
    \item каждый процесс, кроме начального, имеет PPID;
    \item процессы образуют дерево;
    \item системный процесс \texttt{init} или \texttt{systemd} часто является предком многих процессов.
\end{itemize}

\subsection{Дерево процессов}

\begin{center}
\begin{tikzpicture}[
    every node/.style={draw, rounded corners, align=center, minimum width=2.2cm},
    level distance=1.4cm,
    sibling distance=3.5cm
]
\node {systemd\\PID 1}
    child { node {sshd}
        child { node {bash}
            child { node {vim} }
            child { node {gcc} }
        }
    }
    child { node {cron} }
    child { node {dbus-daemon} };
\end{tikzpicture}
\end{center}

\subsection{Модель fork--exec}

В UNIX-подобных ОС создание новой программы часто выполняется в два этапа:

\begin{enumerate}
    \item \texttt{fork()} создает копию текущего процесса.
    \item \texttt{exec()} заменяет образ дочернего процесса новой программой.
\end{enumerate}

После \texttt{fork()}:
\begin{itemize}
    \item родитель и ребенок продолжают выполнение с инструкции после \texttt{fork()};
    \item в родительском процессе \texttt{fork()} возвращает PID ребенка;
    \item в дочернем процессе \texttt{fork()} возвращает $0$;
    \item при ошибке возвращается отрицательное значение.
\end{itemize}

\subsection{Алгоритм создания процесса по модели fork--exec}

\begin{enumerate}
    \item Родительский процесс вызывает \texttt{fork()}.
    \item Ядро создает PCB для нового процесса.
    \item Ядро копирует или логически дублирует адресное пространство родителя.
    \item Создается копия таблицы файловых дескрипторов.
    \item Родителю возвращается PID ребенка.
    \item Ребенку возвращается $0$.
    \item Дочерний процесс вызывает \texttt{exec()}.
    \item Ядро загружает новый исполняемый файл.
    \item Старый код и данные дочернего процесса заменяются новым образом.
    \item Дочерний процесс начинает выполнение новой программы.
    \item Родитель при необходимости вызывает \texttt{wait()} или \texttt{waitpid()}.
\end{enumerate}

\subsection{Минимальный пример fork}

\begin{lstlisting}[language=C]
#include <stdio.h>
#include <unistd.h>
#include <sys/types.h>

int main(void) {
    pid_t pid = fork();

    if (pid < 0) {
        perror("fork");
        return 1;
    }

    if (pid == 0) {
        printf("Child: pid=%d, ppid=%d\n", getpid(), getppid());
    } else {
        printf("Parent: pid=%d, child=%d\n", getpid(), pid);
    }

    return 0;
}
\end{lstlisting}

\subsection{Минимальный пример fork--exec--wait}

\begin{lstlisting}[language=C]
#include <stdio.h>
#include <unistd.h>
#include <sys/wait.h>

int main(void) {
    pid_t pid = fork();

    if (pid == 0) {
        execlp("ls", "ls", "-l", NULL);
        perror("execlp");
        return 1;
    } else if (pid > 0) {
        int status;
        waitpid(pid, &status, 0);
        printf("Child finished\n");
    } else {
        perror("fork");
        return 1;
    }

    return 0;
}
\end{lstlisting}

\subsection{Процессы-сироты и процессы-зомби}

\textbf{Процесс-сирота} --- процесс, родитель которого завершился раньше него. В UNIX-подобных системах такой процесс обычно усыновляется специальным системным процессом.

\textbf{Процесс-зомби} --- завершившийся процесс, запись о котором еще остается в таблице процессов, поскольку родитель не считал его код завершения с помощью \texttt{wait()} или \texttt{waitpid()}.

Зомби-процесс:
\begin{itemize}
    \item уже не выполняется;
    \item не потребляет процессорное время;
    \item занимает запись в таблице процессов;
    \item хранит код завершения до получения его родителем.
\end{itemize}

\section{Состояния процессов}

\subsection{Базовая модель состояний}

Классическая модель процесса включает следующие состояния:

\begin{enumerate}
    \item \textbf{New} --- процесс создается.
    \item \textbf{Ready} --- процесс готов к выполнению и ожидает процессор.
    \item \textbf{Running} --- процесс выполняется на процессоре.
    \item \textbf{Blocked} или \textbf{Waiting} --- процесс ожидает событие.
    \item \textbf{Terminated} --- процесс завершил работу.
\end{enumerate}

\subsection{Диаграмма состояний}

\begin{center}
\begin{tikzpicture}[
    state/.style={draw, rounded corners, minimum width=2.8cm, minimum height=1cm, align=center},
    arrow/.style={->, thick}
]
\node[state] (new) {New};
\node[state, right=2.3cm of new] (ready) {Ready};
\node[state, right=2.3cm of ready] (running) {Running};
\node[state, below=1.6cm of running] (blocked) {Blocked};
\node[state, right=2.3cm of running] (term) {Terminated};

\draw[arrow] (new) -- node[above]{admit} (ready);
\draw[arrow] (ready) -- node[above]{dispatch} (running);
\draw[arrow] (running) -- node[above]{exit} (term);
\draw[arrow] (running) to[bend left=20] node[below]{preemption} (ready);
\draw[arrow] (running) -- node[right]{I/O wait} (blocked);
\draw[arrow] (blocked) -- node[below]{event} (ready);
\end{tikzpicture}
\end{center}

\subsection{Расширенная модель}

В современных ОС могут использоваться дополнительные состояния:

\begin{itemize}
    \item \textbf{Suspended Ready} --- процесс готов, но выгружен из памяти.
    \item \textbf{Suspended Blocked} --- процесс ожидает событие и выгружен из памяти.
    \item \textbf{Stopped} --- процесс остановлен сигналом или отладчиком.
    \item \textbf{Zombie} --- процесс завершен, но его статус еще не получен родителем.
\end{itemize}

\section{Многозадачный и многопрограммный режим}

\subsection{Многопрограммность}

\textbf{Многопрограммность} --- режим работы ОС, при котором в памяти одновременно находится несколько программ, а процессор переключается между ними.

Цель многопрограммности:
\begin{itemize}
    \item повысить загрузку процессора;
    \item скрыть задержки ввода-вывода;
    \item увеличить пропускную способность системы.
\end{itemize}

Если один процесс ожидает диск или сеть, процессор может выполнять другой процесс.

\subsection{Многозадачность}

\textbf{Многозадачность} --- способность ОС обеспечивать видимость одновременного выполнения нескольких задач.

На одноядерной системе физически в каждый момент времени выполняется только один поток, но быстрые переключения создают иллюзию параллельности.

На многоядерной системе возможен реальный параллелизм: несколько потоков действительно выполняются одновременно на разных ядрах.

\subsection{Кооперативная и вытесняющая многозадачность}

\begin{longtable}{|p{4cm}|p{5cm}|p{5cm}|}
\hline
\textbf{Признак} & \textbf{Кооперативная многозадачность} & \textbf{Вытесняющая многозадачность} \\
\hline
Кто отдает процессор & Сам процесс & ОС по таймеру или событию \\
\hline
Устойчивость & Один зависший процесс может заблокировать систему & ОС может принудительно переключить процесс \\
\hline
Реактивность & Зависит от корректности программ & Контролируется планировщиком \\
\hline
Современное применение & Ограниченно & Основной подход в современных ОС \\
\hline
\end{longtable}

\subsection{Квант времени}

\textbf{Квант времени} --- максимальный интервал, в течение которого процесс может непрерывно занимать процессор до вытеснения.

Если квант слишком мал:
\begin{itemize}
    \item увеличиваются накладные расходы на переключение контекста.
\end{itemize}

Если квант слишком велик:
\begin{itemize}
    \item ухудшается интерактивность;
    \item система становится менее отзывчивой.
\end{itemize}

\section{Планирование процессов}

\subsection{Назначение планировщика}

\textbf{Планировщик} --- компонент ОС, выбирающий, какой процесс или поток должен выполняться следующим.

Основные цели планирования:
\begin{itemize}
    \item справедливое распределение процессора;
    \item высокая загрузка CPU;
    \item малое время отклика;
    \item высокая пропускная способность;
    \item соблюдение приоритетов;
    \item поддержка процессов реального времени.
\end{itemize}

\subsection{Очереди планирования}

ОС обычно поддерживает несколько очередей:

\begin{itemize}
    \item очередь новых процессов;
    \item очередь готовых процессов;
    \item очереди ожидания устройств ввода-вывода;
    \item очереди ожидания синхронизационных объектов;
    \item очереди процессов реального времени.
\end{itemize}

\subsection{Критерии эффективности планирования}

\begin{enumerate}
    \item \textbf{Загрузка CPU}
    \[
        U = \frac{T_{\text{busy}}}{T_{\text{total}}}.
    \]

    \item \textbf{Пропускная способность} --- число завершенных процессов в единицу времени.

    \item \textbf{Время оборота}
    \[
        T_{\text{turnaround}} = T_{\text{finish}} - T_{\text{arrival}}.
    \]

    \item \textbf{Время ожидания}
    \[
        T_{\text{waiting}} = T_{\text{turnaround}} - T_{\text{CPU}}.
    \]

    \item \textbf{Время отклика}
    \[
        T_{\text{response}} = T_{\text{first run}} - T_{\text{arrival}}.
    \]
\end{enumerate}

\subsection{Алгоритм FCFS}

\textbf{FCFS} (\textit{First-Come, First-Served}) --- процесс, пришедший первым, обслуживается первым.

\subsubsection*{Шаги алгоритма}

\begin{enumerate}
    \item Все готовые процессы помещаются в FIFO-очередь.
    \item Планировщик выбирает процесс из головы очереди.
    \item Процесс выполняется до завершения или блокировки.
    \item После освобождения процессора выбирается следующий процесс.
\end{enumerate}

\subsubsection*{Пример}

Пусть есть процессы:

\[
\begin{array}{c|c}
\text{Процесс} & \text{CPU burst} \\
\hline
P_1 & 5 \\
P_2 & 3 \\
P_3 & 2
\end{array}
\]

Диаграмма:

\[
\begin{array}{|c|c|c|c|}
\hline
0 & 5 & 8 & 10 \\
\hline
P_1 & P_2 & P_3 & \\
\hline
\end{array}
\]


Время ожидания:
\[
W_1 = 0,\quad W_2 = 5,\quad W_3 = 8.
\]

Среднее время ожидания:
\[
\overline{W} = \frac{0 + 5 + 8}{3} = \frac{13}{3}.
\]

\subsection{Алгоритм Round Robin}

\textbf{Round Robin} --- вытесняющий алгоритм, в котором каждому процессу выделяется квант времени $q$.

\subsubsection*{Шаги алгоритма}

\begin{enumerate}
    \item Все готовые процессы находятся в FIFO-очереди.
    \item Из головы очереди выбирается процесс.
    \item Процесс выполняется не более $q$ единиц времени.
    \item Если процесс завершился раньше, он удаляется из очереди.
    \item Если процесс не завершился, он вытесняется и помещается в конец очереди.
    \item Алгоритм повторяется.
\end{enumerate}

\subsubsection*{Пример}

Пусть:
\[
q = 2,
\]
\[
\begin{array}{c|c}
\text{Процесс} & \text{CPU burst} \\
\hline
P_1 & 5 \\
P_2 & 3 \\
P_3 & 1
\end{array}
\]

Диаграмма:

\[
\begin{array}{|c|c|c|c|c|c|c|}
\hline
0 & 2 & 4 & 5 & 7 & 8 & 9 \\
\hline
P_1 & P_2 & P_3 & P_1 & P_2 & P_1 & \\
\hline
\end{array}
\]

Порядок выполнения:
\[
P_1 \rightarrow P_2 \rightarrow P_3 \rightarrow P_1 \rightarrow P_2 \rightarrow P_1.
\]

\subsection{Приоритетное планирование}

Каждому процессу назначается приоритет. Планировщик выбирает готовый процесс с наибольшим приоритетом.

Проблема: \textbf{голодание} низкоприоритетных процессов.

Решение: \textbf{старение} --- постепенное повышение приоритета процесса, долго ожидающего в очереди.

\subsection{SJF и SRTF}

\textbf{SJF} (\textit{Shortest Job First}) выбирает процесс с минимальной ожидаемой продолжительностью следующего CPU burst.

\textbf{SRTF} (\textit{Shortest Remaining Time First}) --- вытесняющий вариант SJF, выбирающий процесс с минимальным оставшимся временем выполнения.

Преимущество: малое среднее время ожидания.

Недостаток: необходимо оценивать будущую длительность CPU burst.

\section{Управление процессами}

\subsection{Основные операции управления}

ОС предоставляет системные вызовы и команды для:

\begin{itemize}
    \item создания процесса;
    \item завершения процесса;
    \item ожидания завершения;
    \item изменения приоритета;
    \item отправки сигналов;
    \item остановки и продолжения;
    \item получения информации;
    \item трассировки и отладки.
\end{itemize}

\subsection{Типичные системные вызовы UNIX-подобных ОС}

\begin{longtable}{|p{4cm}|p{9cm}|}
\hline
\textbf{Системный вызов} & \textbf{Назначение} \\
\hline
\texttt{fork()} & Создать дочерний процесс \\
\hline
\texttt{exec()} & Заменить образ процесса новой программой \\
\hline
\texttt{exit()} & Завершить процесс \\
\hline
\texttt{wait()} & Дождаться завершения дочернего процесса \\
\hline
\texttt{waitpid()} & Дождаться конкретного дочернего процесса \\
\hline
\texttt{getpid()} & Получить PID текущего процесса \\
\hline
\texttt{getppid()} & Получить PID родителя \\
\hline
\texttt{kill()} & Отправить сигнал процессу \\
\hline
\texttt{nice()} & Изменить относительный приоритет \\
\hline
\texttt{setpriority()} & Установить приоритет процесса \\
\hline
\end{longtable}

\subsection{Команды управления процессами}

\begin{longtable}{|p{4cm}|p{9cm}|}
\hline
\textbf{Команда} & \textbf{Назначение} \\
\hline
\texttt{ps} & Показать список процессов \\
\hline
\texttt{top} & Интерактивно показать процессы и загрузку системы \\
\hline
\texttt{htop} & Расширенный интерактивный монитор процессов \\
\hline
\texttt{pgrep} & Найти PID по имени или свойствам процесса \\
\hline
\texttt{pidof} & Найти PID программы \\
\hline
\texttt{kill} & Отправить сигнал процессу \\
\hline
\texttt{pkill} & Отправить сигнал процессам по имени или шаблону \\
\hline
\texttt{killall} & Завершить процессы по имени \\
\hline
\texttt{jobs} & Показать задания текущей оболочки \\
\hline
\texttt{fg} & Перевести задание на передний план \\
\hline
\texttt{bg} & Продолжить задание в фоне \\
\hline
\texttt{nice} & Запустить процесс с измененным приоритетом \\
\hline
\texttt{renice} & Изменить приоритет уже запущенного процесса \\
\hline
\texttt{pstree} & Показать дерево процессов \\
\hline
\texttt{strace} & Отслеживать системные вызовы процесса \\
\hline
\texttt{lsof} & Показать открытые файлы процесса \\
\hline
\end{longtable}

\subsection{Примеры команд}

Показать процессы текущего пользователя:

\begin{lstlisting}[language=bash]
ps u
\end{lstlisting}

Показать все процессы:

\begin{lstlisting}[language=bash]
ps aux
\end{lstlisting}

Найти процесс по имени:

\begin{lstlisting}[language=bash]
pgrep nginx
\end{lstlisting}

Завершить процесс по PID:

\begin{lstlisting}[language=bash]
kill 1234
\end{lstlisting}

Принудительно завершить процесс:

\begin{lstlisting}[language=bash]
kill -9 1234
\end{lstlisting}

Остановить процесс:

\begin{lstlisting}[language=bash]
kill -STOP 1234
\end{lstlisting}

Продолжить процесс:

\begin{lstlisting}[language=bash]
kill -CONT 1234
\end{lstlisting}

Запустить программу в фоне:

\begin{lstlisting}[language=bash]
long_task &
\end{lstlisting}

Посмотреть фоновые задания:

\begin{lstlisting}[language=bash]
jobs
\end{lstlisting}

Вернуть задание на передний план:

\begin{lstlisting}[language=bash]
fg %1
\end{lstlisting}

\section{Сигналы}

\subsection{Понятие сигнала}

\textbf{Сигнал} --- асинхронное уведомление процессу о некотором событии.

Сигналы используются для:
\begin{itemize}
    \item завершения процесса;
    \item остановки процесса;
    \item продолжения процесса;
    \item уведомления о завершении дочернего процесса;
    \item обработки ошибок;
    \item пользовательского взаимодействия.
\end{itemize}

\subsection{Основные сигналы UNIX}

\begin{longtable}{|p{3cm}|p{3cm}|p{7cm}|}
\hline
\textbf{Сигнал} & \textbf{Номер обычно} & \textbf{Назначение} \\
\hline
\texttt{SIGHUP} & 1 & Потеря управляющего терминала или перечитывание конфигурации \\
\hline
\texttt{SIGINT} & 2 & Прерывание с клавиатуры, обычно \texttt{Ctrl+C} \\
\hline
\texttt{SIGKILL} & 9 & Безусловное завершение, нельзя перехватить \\
\hline
\texttt{SIGTERM} & 15 & Вежливая просьба завершиться \\
\hline
\texttt{SIGSTOP} & обычно 19 & Безусловная остановка, нельзя перехватить \\
\hline
\texttt{SIGCONT} & обычно 18 & Продолжить остановленный процесс \\
\hline
\texttt{SIGCHLD} & обычно 17 & Завершился или изменил состояние дочерний процесс \\
\hline
\end{longtable}

\subsection{Обработка сигналов}

Для большинства сигналов процесс может:

\begin{itemize}
    \item выполнить действие по умолчанию;
    \item проигнорировать сигнал;
    \item установить собственный обработчик.
\end{itemize}

Исключения:
\begin{itemize}
    \item \texttt{SIGKILL} нельзя перехватить или игнорировать;
    \item \texttt{SIGSTOP} нельзя перехватить или игнорировать.
\end{itemize}

\subsection{Пример обработки SIGINT}

\begin{lstlisting}[language=C]
#include <stdio.h>
#include <signal.h>
#include <unistd.h>

void handler(int signo) {
    printf("Caught signal %d\n", signo);
}

int main(void) {
    signal(SIGINT, handler);

    while (1) {
        printf("Working...\n");
        sleep(1);
    }

    return 0;
}
\end{lstlisting}

\section{Взаимодействие процессов}

\subsection{Независимые и взаимодействующие процессы}

\textbf{Независимый процесс} не влияет на другие процессы и не зависит от них.

\textbf{Взаимодействующие процессы} обмениваются данными или синхронизируют действия.

Причины взаимодействия:
\begin{itemize}
    \item обмен данными;
    \item ускорение вычислений;
    \item модульность;
    \item совместное использование ресурсов;
    \item построение клиент-серверных приложений.
\end{itemize}

\subsection{Основные модели IPC}

\textbf{IPC} (\textit{Inter-Process Communication}) --- межпроцессное взаимодействие.

Две базовые модели:

\begin{enumerate}
    \item \textbf{Разделяемая память}.

    Процессы получают доступ к общей области памяти.

    \[
        P_1 \leftrightarrow \text{shared memory} \leftrightarrow P_2
    \]

    Преимущества:
    \begin{itemize}
        \item высокая скорость;
        \item нет необходимости копировать данные через ядро после настройки.
    \end{itemize}

    Недостатки:
    \begin{itemize}
        \item требуется синхронизация;
        \item возможны гонки данных;
        \item сложнее обеспечить безопасность.
    \end{itemize}

    \item \textbf{Передача сообщений}.

    Процессы обмениваются сообщениями через ядро.

    \[
        P_1 \rightarrow \text{kernel} \rightarrow P_2
    \]

    Преимущества:
    \begin{itemize}
        \item проще синхронизация;
        \item хорошая изоляция;
        \item удобно для распределенных систем.
    \end{itemize}

    Недостатки:
    \begin{itemize}
        \item возможны дополнительные копирования;
        \item выше накладные расходы.
    \end{itemize}
\end{enumerate}

\subsection{Средства IPC}

\begin{longtable}{|p{4cm}|p{9cm}|}
\hline
\textbf{Средство} & \textbf{Описание} \\
\hline
Каналы \texttt{pipe} & Однонаправленный поток байтов между родственными процессами \\
\hline
Именованные каналы FIFO & Канал с именем в файловой системе, может связывать неродственные процессы \\
\hline
Очереди сообщений & Передача дискретных сообщений через ядро \\
\hline
Разделяемая память & Общая область памяти для нескольких процессов \\
\hline
Семафоры & Синхронизация доступа к общим ресурсам \\
\hline
Мьютексы & Взаимное исключение, чаще внутри процесса или в разделяемой памяти \\
\hline
Условные переменные & Ожидание наступления условия \\
\hline
Сигналы & Асинхронные уведомления \\
\hline
Сокеты & Локальное или сетевое взаимодействие процессов \\
\hline
Файлы & Простейшая форма обмена данными через файловую систему \\
\hline
Отображение файлов в память & Совместное использование данных через \texttt{mmap} \\
\hline
\end{longtable}

\section{Каналы}

\subsection{Неименованный канал}

\textbf{Канал} --- однонаправленный буфер, обычно работающий по принципу FIFO.

В UNIX канал создается системным вызовом:

\begin{lstlisting}[language=C]
int pipe(int fd[2]);
\end{lstlisting}

После вызова:
\begin{itemize}
    \item \texttt{fd[0]} --- конец для чтения;
    \item \texttt{fd[1]} --- конец для записи.
\end{itemize}

\subsection{Алгоритм использования pipe с fork}

\begin{enumerate}
    \item Родитель вызывает \texttt{pipe(fd)}.
    \item Родитель вызывает \texttt{fork()}.
    \item Дочерний процесс закрывает ненужный конец канала.
    \item Родительский процесс закрывает ненужный конец канала.
    \item Один процесс пишет в \texttt{fd[1]}.
    \item Другой процесс читает из \texttt{fd[0]}.
    \item После завершения обмена оба процесса закрывают дескрипторы.
\end{enumerate}

\subsection{Пример pipe}

\begin{lstlisting}[language=C]
#include <stdio.h>
#include <unistd.h>
#include <string.h>
#include <sys/wait.h>

int main(void) {
    int fd[2];
    char buf[100];

    if (pipe(fd) == -1) {
        perror("pipe");
        return 1;
    }

    pid_t pid = fork();

    if (pid == 0) {
        close(fd[0]);
        const char *msg = "Hello from child";
        write(fd[1], msg, strlen(msg) + 1);
        close(fd[1]);
    } else {
        close(fd[1]);
        read(fd[0], buf, sizeof(buf));
        printf("Parent received: %s\n", buf);
        close(fd[0]);
        wait(NULL);
    }

    return 0;
}
\end{lstlisting}

\subsection{Командные конвейеры}

Командная оболочка использует каналы для построения конвейеров:

\begin{lstlisting}[language=bash]
ps aux | grep nginx | wc -l
\end{lstlisting}

Смысл:
\begin{enumerate}
    \item \texttt{ps aux} выводит список процессов.
    \item \texttt{grep nginx} фильтрует строки.
    \item \texttt{wc -l} считает количество строк.
\end{enumerate}

\section{Разделяемая память}

\subsection{Общее понятие}

\textbf{Разделяемая память} --- область памяти, отображенная в адресные пространства нескольких процессов.

Схема:

\begin{center}
\begin{tikzpicture}[
    block/.style={draw, rounded corners, minimum width=3cm, minimum height=1cm, align=center},
    arrow/.style={<->, thick}
]
\node[block] (p1) {Процесс $P_1$};
\node[block, right=2cm of p1] (mem) {Shared memory};
\node[block, right=2cm of mem] (p2) {Процесс $P_2$};

\draw[arrow] (p1) -- (mem);
\draw[arrow] (mem) -- (p2);
\end{tikzpicture}
\end{center}

\subsection{Проблема синхронизации}

Если два процесса одновременно изменяют общие данные, возникает \textbf{гонка данных}.

Пример опасной операции:

\[
    x := x + 1.
\]

На машинном уровне она может состоять из нескольких шагов:

\begin{enumerate}
    \item загрузить $x$ в регистр;
    \item увеличить регистр;
    \item записать регистр обратно в память.
\end{enumerate}

Если два процесса выполняют эти действия одновременно, одно обновление может потеряться.

\section{Критические секции и синхронизация}

\subsection{Критическая секция}

\textbf{Критическая секция} --- участок кода, в котором процесс обращается к общему ресурсу.

Задача синхронизации --- организовать выполнение так, чтобы несколько процессов не разрушали общие данные.

\subsection{Требования к решению задачи критической секции}

Корректное решение должно удовлетворять трем условиям:

\begin{enumerate}
    \item \textbf{Взаимное исключение}.

    Если один процесс находится в критической секции, другой процесс не может одновременно находиться в своей критической секции для того же ресурса.

    \item \textbf{Прогресс}.

    Если ни один процесс не находится в критической секции, а некоторые хотят войти, выбор следующего входящего не должен откладываться бесконечно.

    \item \textbf{Ограниченное ожидание}.

    Существует конечная верхняя граница на число раз, которое другие процессы могут войти в критическую секцию после того, как данный процесс запросил вход.
\end{enumerate}

\subsection{Семафор}

\textbf{Семафор} --- целочисленный синхронизационный объект, над которым определены две атомарные операции:

\begin{itemize}
    \item \texttt{wait}, также называемая $P$ или \texttt{down};
    \item \texttt{signal}, также называемая $V$ или \texttt{up}.
\end{itemize}

Операции:

\[
\texttt{wait}(S):
\begin{cases}
S := S - 1, & \text{если } S > 0,\\
\text{заблокировать процесс}, & \text{иначе.}
\end{cases}
\]

\[
\texttt{signal}(S):
\begin{cases}
\text{разбудить ожидающий процесс}, & \text{если такие есть},\\
S := S + 1, & \text{иначе.}
\end{cases}
\]

\subsection{Бинарный семафор}

Бинарный семафор принимает значения $0$ и $1$ и используется как мьютекс.

\begin{lstlisting}
wait(mutex);
critical_section();
signal(mutex);
\end{lstlisting}

\subsection{Мьютекс}

\textbf{Мьютекс} --- объект взаимного исключения.

Типичная схема:

\begin{lstlisting}
lock(mutex);
critical_section();
unlock(mutex);
\end{lstlisting}

В отличие от общего семафора, мьютекс часто имеет владельца: разблокировать его должен тот поток или процесс, который его захватил.

\section{Классическая задача производитель--потребитель}

\subsection{Постановка}

Есть:
\begin{itemize}
    \item производитель, помещающий элементы в буфер;
    \item потребитель, извлекающий элементы из буфера;
    \item буфер размера $N$.
\end{itemize}

Необходимо:
\begin{itemize}
    \item не писать в полный буфер;
    \item не читать из пустого буфера;
    \item не допускать одновременного разрушительного доступа к буферу.
\end{itemize}

\subsection{Решение с семафорами}

Используются:
\begin{itemize}
    \item \texttt{mutex = 1} --- взаимное исключение;
    \item \texttt{empty = N} --- число свободных ячеек;
    \item \texttt{full = 0} --- число занятых ячеек.
\end{itemize}

Производитель:

\begin{lstlisting}
while true:
    item = produce()
    wait(empty)
    wait(mutex)
    insert(item)
    signal(mutex)
    signal(full)
\end{lstlisting}

Потребитель:

\begin{lstlisting}
while true:
    wait(full)
    wait(mutex)
    item = remove()
    signal(mutex)
    signal(empty)
    consume(item)
\end{lstlisting}

\subsection{Пояснение шагов}

\begin{enumerate}
    \item Производитель сначала ожидает свободную ячейку.
    \item Затем входит в критическую секцию.
    \item Помещает элемент в буфер.
    \item Выходит из критической секции.
    \item Увеличивает счетчик занятых ячеек.
    \item Потребитель сначала ожидает наличие элемента.
    \item Затем входит в критическую секцию.
    \item Извлекает элемент.
    \item Выходит из критической секции.
    \item Увеличивает счетчик свободных ячеек.
\end{enumerate}

\section{Взаимные блокировки}

\subsection{Понятие deadlock}

\textbf{Взаимная блокировка} --- ситуация, в которой несколько процессов бесконечно ожидают события, которое может быть вызвано только одним из ожидающих процессов.

Пример:

\begin{itemize}
    \item $P_1$ захватил ресурс $R_1$ и ждет $R_2$;
    \item $P_2$ захватил ресурс $R_2$ и ждет $R_1$.
\end{itemize}

\subsection{Условия Коффмана}

Для возникновения deadlock необходимы четыре условия:

\begin{enumerate}
    \item \textbf{Взаимное исключение}.

    Некоторые ресурсы не могут использоваться несколькими процессами одновременно.

    \item \textbf{Удержание и ожидание}.

    Процесс удерживает один ресурс и ожидает другой.

    \item \textbf{Отсутствие принудительного изъятия}.

    Ресурс не может быть отобран у процесса насильно.

    \item \textbf{Циклическое ожидание}.

    Существует цикл процессов, каждый из которых ожидает ресурс, удерживаемый следующим процессом в цикле.
\end{enumerate}

\subsection{Методы борьбы с deadlock}

\begin{enumerate}
    \item \textbf{Предотвращение}.

    Исключить хотя бы одно из условий Коффмана.

    \item \textbf{Избежание}.

    Выдавать ресурсы только если система остается в безопасном состоянии.

    \item \textbf{Обнаружение и восстановление}.

    Разрешать deadlock возникать, затем обнаруживать его и устранять.

    \item \textbf{Игнорирование}.

    Если deadlock редок, ОС может не заниматься его автоматическим устранением.
\end{enumerate}

\section{Защита и изоляция процессов}

\subsection{Адресное пространство}

Современные ОС обычно предоставляют каждому процессу собственное виртуальное адресное пространство.

Это дает:
\begin{itemize}
    \item изоляцию процессов;
    \item защиту памяти;
    \item удобство размещения программ;
    \item возможность использовать виртуальную память;
    \item независимость адресов разных процессов.
\end{itemize}

\subsection{Режим пользователя и режим ядра}

Процессоры поддерживают разные уровни привилегий.

\begin{itemize}
    \item \textbf{Режим пользователя} --- обычный код приложений, ограниченный доступ к памяти и устройствам.
    \item \textbf{Режим ядра} --- привилегированный режим, в котором выполняется ядро ОС.
\end{itemize}

Переход в ядро происходит через:
\begin{itemize}
    \item системные вызовы;
    \item исключения;
    \item аппаратные прерывания.
\end{itemize}

\section{Процессы в современных ОС}

\subsection{UNIX/Linux}

Характерные черты:
\begin{itemize}
    \item древовидная структура процессов;
    \item модель \texttt{fork}/\texttt{exec};
    \item файловые дескрипторы наследуются потомками;
    \item сигналы как механизм уведомления;
    \item множество IPC-средств;
    \item процессы и потоки в ядре представлены близкими сущностями.
\end{itemize}

\subsection{Windows}

В Windows процесс является контейнером ресурсов, а поток --- единицей выполнения.

Процесс Windows обычно содержит:
\begin{itemize}
    \item виртуальное адресное пространство;
    \item дескрипторы объектов ядра;
    \item маркер доступа;
    \item один или несколько потоков;
    \item информацию о приоритетах.
\end{itemize}

Для IPC применяются:
\begin{itemize}
    \item каналы;
    \item именованные каналы;
    \item сокеты;
    \item разделяемая память;
    \item mailslots;
    \item RPC;
    \item COM;
    \item события, мьютексы, семафоры.
\end{itemize}

\section{Краткое резюме}

Процесс --- это выполняющаяся программа вместе с ее ресурсами и состоянием. ОС представляет процессы с помощью специальных управляющих структур, содержащих идентификаторы, состояние, контекст, сведения о памяти, файлах и планировании.

Процессы могут находиться в состояниях создания, готовности, выполнения, ожидания и завершения. Многозадачность достигается за счет планирования и переключения контекста. Вытесняющая многозадачность позволяет ОС самостоятельно отбирать процессор у процесса по таймеру или приоритету.

Процессы образуют иерархии порождения. В UNIX-подобных ОС распространена модель \texttt{fork--exec--wait}. Управление процессами осуществляется системными вызовами и пользовательскими командами, такими как \texttt{ps}, \texttt{kill}, \texttt{jobs}, \texttt{fg}, \texttt{bg}, \texttt{nice}, \texttt{renice}.

Межпроцессное взаимодействие необходимо для обмена данными и синхронизации. Основные подходы --- разделяемая память и передача сообщений. Практические средства IPC включают каналы, FIFO, очереди сообщений, сокеты, сигналы, семафоры, мьютексы и разделяемую память.

\end{document}
